% !TeX root = ../Thesis.tex

%*******************************************************
% Abstract
%*******************************************************
%\renewcommand{\abstractname}{Abstract}
\pdfbookmark[0]{Abstract}{Abstract}
% \addcontentsline{toc}{chapter}{\tocEntry{Abstract}}
\begingroup
\let\clearpage\relax
\let\cleardoublepage\relax
\let\cleardoublepage\relax

\chapter*{Abstract}
%\myAbstract{}
[WORK IN PROGRESS] Sharded distributed databases have gained significant attention due to the limitations of traditional databases in handling large-scale data. These limitations of traditional databases in handling large-scale data are addressed through sharding, which distributes data across multiple database nodes, allowing for better performance, increased scalability, and improved fault tolerance. This study conducts an analysis of existing literature and implements a prototype system using a middleware approach to enable a system of distributed databases. The prototype system is evaluated using custom benchmarking techniques, considering factors such as performance and data management. The results show that sharding can effectively scale databases horizontally, but its implementation requires careful design and management to ensure optimal performance and data consistency. The impact of this research and implementation is demonstrating the effectiveness of consistent hash sharding in Redis key-value stores as a scalable solution for large-scale data handling.
%Short summary of the contents in English\dots a great guide by
%Kent Beck how to write good abstracts can be found here:
%\begin{center}
%\url{https://plg.uwaterloo.ca/~migod/research/beckOOPSLA.html}
%\end{center}
%Steve Easterbrook also has a good article on the same topic:
%\begin{center}
%\url{https://www.easterbrook.ca/steve/2010/01/how-to-write-a-scientific-abstract-in-six-easy-steps/}
%\end{center}

\vfill

\begin{otherlanguage}{ngerman}
\pdfbookmark[0]{Zusammenfassung}{Zusammenfassung}
\chapter*{Zusammenfassung}
%Kurze Zusammenfassung des Inhaltes in deutscher Sprache\dots
\end{otherlanguage}

\endgroup

\vfill
